\clearpage
\section*{Problem 1}
\addcontentsline{toc}{section}{Problem 1}
\begingroup
\hypersetup{colorlinks=true,
            linkcolor=blue!50!black,
            citecolor=blue!50!black,
            urlcolor=blue!50!black}
\newcommand{\qoneT}{\mathbb{T}}
\newcommand{\qoneR}{\mathbb{R}}
\newcommand{\qoneN}{\mathbb{N}}
\newcommand{\qoneC}{\mathcal{C}}
\newcommand{\qoneD}{\mathcal{D}}
\newcommand{\qoneE}{\mathbb{E}}
\newcommand{\qoneHp}{H}      % Hermite
\newcommand{\qonePN}{P_N}
\newcommand{\qoneLin}{\mathbf{l}}   % stationary linear solution (placeholder for tree symbol)
\newcommand{\qoneLsq}{\mathbf{q}}   % Wick square
\newcommand{\qoneLcb}{\mathbf{c}}   % Wick cube
\newcommand{\qoneIqc}{\mathbf{C}}   % solution driven by the cube
\newtheorem{qonetheorem}{Theorem}
\newtheorem{qonelemma}{Lemma}
\newtheorem{qoneproposition}{Proposition}
\newtheorem{qonecorollary}{Corollary}
\newtheorem{remark}{Remark}
\title{Solution to Problem 1: (Lack of) Quasi-Shift-Invariance\\
of the $\Phi^4_3$ Measure}
\author{}
\date{}
\maketitle


\section*{Statement and Answer}

Let $\qoneT^3$ be the unit $3$-torus and $\mu$ the $\Phi^4_3$ measure on
$\qoneD'(\qoneT^3)$. For smooth $\psi\not\equiv0$ let $T_\psi(u)=u+\psi$, and let
$T_\psi^*\mu$ be the pushforward.

\medskip
\noindent\textbf{Theorem.} \emph{For every smooth $\psi\not\equiv0$ and every
admissible choice of mass and (nonzero) coupling constant, the measures $\mu$ and
$T_\psi^*\mu$ are mutually \textbf{singular}.}

\medskip
\noindent The answer is \textbf{NO}: $\mu$ and $T_\psi^*\mu$ are \emph{not}
equivalent. This is a simplified form of a result of Hairer--Peroni and rests on
ideas of Hairer--Kusuoka--Nagoji.

\medskip
\noindent\textbf{Warning on a fatal false premise.} It is tempting to write
$d\mu/d\mu_0=Z^{-1}e^{-\mathcal R(u)}$ against the Gaussian free field $\mu_0$
and deduce quasi-invariance by a Cameron--Martin/Wick-binomial computation. This
is \emph{wrong in three dimensions}: Glimm showed that, unlike $d=2$, the
$\Phi^4_3$ measure $\mu$ and the free field are mutually \emph{singular} in
$d=3$ (the borderline dimension is $8/3$). There is no such Radon--Nikodym
density, so any argument built on it collapses. The correct proof exhibits an
event separating $\mu$ from $T_\psi^*\mu$, driven by a genuinely
$3$-dimensional logarithmic divergence.

\section*{Notation}

Fix a space--time white noise $\xi$ on $\qoneR\times\qoneT^3$ and let $\qoneLin$ be the
stationary solution of $(\partial_t+1-\Delta)\qoneLin=\xi$. Write $\qonePN=P_{\le N}$ for
the projection onto Fourier modes $|k|\le N$, $\qoneLin_N=\qonePN\qoneLin$, and
$c_N=\qoneE\qoneLin_N^2$. Let $\qoneHp_n$ be the Hermite polynomials normalized by
$\qoneHp_0\equiv1$, $\qoneHp_n'=n\qoneHp_{n-1}$, $\qoneE\qoneHp_n(Z,1)=0$ for standard normal $Z$. Set
\[
\qoneLsq=\lim_{N}\qoneHp_2(\qoneLin_N,c_N),\qquad \qoneLcb=\lim_N\qoneHp_3(\qoneLin_N,c_N),
\]
the Wick square and cube (convergent in $\qoneC^{-1-2\kappa}$ and $\qoneC^{-3/2-3\kappa}$
respectively), and let $\qoneIqc$ solve $(\partial_t+1-\Delta)\qoneIqc=\qoneLcb$. We fix
$0<\kappa\le 1/100$. By \cite{HKN24}, $\qoneLcb$ and $\qoneIqc$ are a.s.\ continuous in
time with values in $\qoneC^{-1/2-\kappa}$ and $\qoneC^{1/2-3\kappa}$.

\paragraph{Wick chaos and stationary kernels.}
We write $\qoneLin_N(x)=\sum_{|k|\le N}\hat\qoneLin(k)e^{ik\cdot x}$ with
$\qoneE[\hat\qoneLin(k)\hat\qoneLin(-k)]=s(k)$ the equal-time variance of the stationary
Ornstein--Uhlenbeck mode, so that $s(k)=\tfrac12\langle k\rangle^{-2}$,
$\langle k\rangle^2=1+|k|^2$, and $c_N=\sum_{|k|\le N}s(k)$. The Wick powers are the
chaos elements
\[
\qoneLsq_N(x)=\!\!\sum_{|k_1|,|k_2|\le N}\!\!\hat\qoneLin(k_1)\hat\qoneLin(k_2)\,e^{i(k_1+k_2)\cdot x},\quad
\qoneLcb_N(x)=\!\!\sum_{|k_1|,|k_2|,|k_3|\le N}\!\!\hat\qoneLin(k_1)\hat\qoneLin(k_2)\hat\qoneLin(k_3)\,e^{i(k_1+k_2+k_3)\cdot x},
\]
i.e.\ $\qoneLsq_N=\qoneHp_2(\qoneLin_N,c_N)$ lives in the second homogeneous Wick chaos and
$\qoneLcb_N=\qoneHp_3(\qoneLin_N,c_N)$ in the third. The process $\qoneIqc$ is the
stationary solution of the linear heat equation driven by $\qoneLcb$; it is the
image of $\qoneLcb$ under a deterministic Fourier multiplier and therefore
\emph{lives in the third Wick chaos as well}. Its equal-time kernel is
\begin{equation}\label{eq:Ckernel}
\qoneIqc_N(x)=\!\!\sum_{|p_1|,|p_2|,|p_3|\le N}\!\!K(p_1,p_2,p_3)\,
\hat\qoneLin(p_1)\hat\qoneLin(p_2)\hat\qoneLin(p_3)\,e^{i(p_1+p_2+p_3)\cdot x},
\end{equation}
where $K$ is the symmetric Duhamel multiplier $K(p_1,p_2,p_3)=\bigl(\langle p_1\rangle^2+\langle p_2\rangle^2+\langle p_3\rangle^2\bigr)^{-1}$ at equal time (its precise form is irrelevant below; only its symmetry and the bound $0<K\le\langle p_1+p_2+p_3\rangle^{-2}+\dots$ are used).

\paragraph{The resonance constant $c_{N,2}$.}
We first record \emph{why} the naive definition is wrong, and then give the
correct one. Because $\qoneLsq_N$ lies in the second homogeneous Wick chaos and
$\qoneIqc_N$ in the third, orthogonality of the Wick chaoses gives
\begin{equation}\label{eq:vanish}
\qoneE\bigl[\qoneLsq_N(x)\,\qoneIqc_N(x)\bigr]=0\qquad\text{for every }x ,
\end{equation}
\emph{identically} in $N$ (equivalently: under the reflection $\qoneLin\mapsto-\qoneLin$,
which preserves the law of $\qoneLin$, the field $\qoneLsq_N=\qoneHp_2(\qoneLin_N,c_N)$ is even
while $\qoneIqc_N$ is odd, so the product is odd and has zero mean). Hence the
\emph{full expectation} $\qoneE[\qoneLsq_N\qoneIqc_N]$ is \textbf{not} the
renormalization constant; defining $c_{N,2}$ by it would make it vacuously $0$.

The object that genuinely diverges is the \emph{local (diagonal) self-contraction}
hidden inside the pointwise product $\qoneLsq_N\,\qoneIqc_N$. Concretely, the product of a
second- and a third-chaos field decomposes into orthogonal chaos components by the
Wick product formula
\begin{equation}\label{eq:wickprod}
\qoneLsq_N\,\qoneIqc_N=[\![5]\!]_N+[\![3]\!]_N+[\![1]\!]_N ,
\end{equation}
where $[\![5-2k]\!]_N$ is the order-$(5-2k)$ chaos component obtained by Wick-contracting
$k$ legs of $\qoneLsq_N$ with $k$ legs of $\qoneIqc_N$ ($k=0,1,2$); cf.\ \cite{CC18}. Each
contraction of a leg of frequency $q$ of $\qoneLsq_N$ with a leg of frequency $-q$ of
$\qoneIqc_N$ contributes the variance factor $s(q)=\qoneE[\hat\qoneLin(q)\hat\qoneLin(-q)]$ together with
the Duhamel weight $K$ of \eqref{eq:Ckernel}; this is a partial contraction, not the
\emph{full} contraction \eqref{eq:vanish}, and is therefore not constrained to vanish.

The divergent counterterm comes from the \emph{doubly contracted} component
$[\![1]\!]_N$, in which both legs of $\qoneLsq_N$ are contracted against two legs of
$\qoneIqc_N$, leaving a single surviving (first-chaos) leg of frequency $m$. Its
combinatorial weight is $2!\binom22\binom32=6$, and carrying out the two internal
contractions gives the field
\begin{equation}\label{eq:order1field}
[\![1]\!]_N(x)=6\!\!\sum_{|m|\le N}\Bigl(\sum_{|q_1|,|q_2|\le N}s(q_1)\,s(q_2)\,K(-q_1,-q_2,m)\Bigr)\hat\qoneLin(m)\,e^{im\cdot x}.
\end{equation}
The inner sum has a frequency-independent leading part that diverges; we extract it
as the renormalization constant
\begin{equation}\label{eq:cN2}
\boxed{\;c_{N,2}\;:=\;6\sum_{|q_1|,|q_2|\le N}s(q_1)\,s(q_2)\,K(-q_1,-q_2,0)
\;=\;3\!\!\sum_{|q_1|,|q_2|\le N}\frac{1}{\langle q_1\rangle^2\langle q_2\rangle^2\bigl(\langle q_1\rangle^2+\langle q_2\rangle^2+1\bigr)}\;}
\end{equation}
(using $s(q)=\tfrac12\langle q\rangle^{-2}$ and $K(-q_1,-q_2,0)=(\langle q_1\rangle^2+\langle q_2\rangle^2+1)^{-1}$).
This is a manifestly \emph{positive}, explicit deterministic constant; it is
\textbf{not} an expectation of $\qoneLsq_N\qoneIqc_N$, and in particular escapes the
vanishing \eqref{eq:vanish}. Lemma~\ref{lem:logdiv} shows $c_{N,2}\simeq\log N$, so the
renormalization is genuinely nontrivial.

\begin{remark}\label{rem:masstype}
The remaining frequency-dependent part of \eqref{eq:order1field}, namely
$[\![1]\!]_N-c_{N,2}\qoneLin_N$, has an \emph{absolutely summable} kernel
(the difference $K(-q_1,-q_2,m)-K(-q_1,-q_2,0)=O(\langle m\rangle^2/(\langle q_1\rangle^2+\langle q_2\rangle^2)^2)$
gains two powers of decay in the loop momenta), hence converges as $N\to\infty$ with
no further renormalization. Likewise the higher components $[\![5]\!]_N$ (no
contraction) and $[\![3]\!]_N$ (a single contraction, whose surviving kernel
$\sum_q s(q)K(-q,d,e)$ is summable in the loop momentum $q$ in three dimensions and
thus finite per output mode) carry no $\log N$ divergence. Thus the \emph{only}
counterterm needed to renormalize $\qoneLsq_N\circ\qoneIqc_N$ is the mass-type counterterm
$c_{N,2}\,\qoneLin_N$, proportional to the first-chaos field $\qoneLin_N$; this is the precise
statement proved in Lemma~\ref{lem:std}. (Replacing $\qoneIqc$ by its frequency truncation
$\qoneIqc_N$ throughout makes every sum finite and \eqref{eq:cN2} unambiguous.)
\end{remark}

\section*{The decomposition of $\mu$}

The statement we use about $\mu$ has two parts: a \emph{static} part, which is the
only thing actually fed into the distinguishing-event computation, and a
\emph{dynamic} part, which is the vehicle by which the static part is proved. We
first state the static fact we need (Proposition~\ref{prop:decomp}), then prove it
by descending from the stationary $\Phi^4_3$ dynamics, addressing explicitly the
two points raised in the construction: (i) that the fixed-time marginal of the
stationary dynamics is exactly $\mu$, and (ii) that the pathwise paracontrolled
structure descends to a $\mu$-almost-everywhere property of the static field $u$.

\begin{qoneproposition}[Static decomposition of a $\mu$-typical field]\label{prop:decomp}
Let $\mu$ be the $\Phi^4_3$ measure on $\qoneD'(\qoneT^3)$. There exists a
probability space $(\Omega,\mathcal F,\mathbb P)$ carrying the stationary linear
solution $\qoneLin$ and its polynomials $\qoneLsq,\qoneLcb,\qoneIqc$ (the
``stochastic data''), and a measurable map
\[
\Psi:\ \omega\longmapsto \bigl(\qoneLin(0,\cdot)(\omega),\,\qoneIqc(0,\cdot)(\omega),\,v(0,\cdot)(\omega),\,v^\sharp(0,\cdot)(\omega)\bigr),
\]
with the following properties.
\begin{enumerate}
\item[\rm(a)] \emph{(Law.)} The field
$u:=\qoneLin(0,\cdot)-\qoneIqc(0,\cdot)+v(0,\cdot)$ has law $\mu$ on
$\qoneD'(\qoneT^3)$; i.e.\ $\mathrm{Law}_{\mathbb P}(u)=\mu$.
\item[\rm(b)] \emph{(Regularity.)} $\mathbb P$-a.s.,
$\qoneLin(0,\cdot)\in\qoneC^{-1/2-\kappa}$,
$\qoneLsq(0,\cdot)\in\qoneC^{-1-2\kappa}$,
$\qoneIqc(0,\cdot)\in\qoneC^{1/2-3\kappa}$,
$v(0,\cdot)\in\qoneC^{1-2\kappa}$, and $v^\sharp(0,\cdot)\in\qoneC^{1+4\kappa}$.
\item[\rm(c)] \emph{(Paracontrolled identity.)} $\mathbb P$-a.s., the static fields
satisfy
\begin{equation}\label{eq:v}
v(0,\cdot)=-3\bigl((v(0,\cdot)-\qoneIqc(0,\cdot))\prec\qoneLin(0,\cdot)\bigr)+v^\sharp(0,\cdot).
\end{equation}
\end{enumerate}
Consequently there is a Borel set $G\subseteq\qoneD'(\qoneT^3)$ with $\mu(G)=1$ such
that every $u\in G$ admits a representation
\begin{equation}\label{eq:u}
u=\ell-C+w,\qquad \ell\in\qoneC^{-1/2-\kappa},\ C\in\qoneC^{1/2-3\kappa},\ w\in\qoneC^{1-2\kappa},
\end{equation}
together with the data $\ell^{\hspace{0.5pt}2}:=\qoneHp_2(\ell)\in\qoneC^{-1-2\kappa}$
and $w^\sharp\in\qoneC^{1+4\kappa}$ obeying
$w=-3((w-C)\prec\ell)+w^\sharp$, where $(\ell,\ell^{\hspace{0.5pt}2},C)$ is in the
closure of the set of admissible smooth stochastic data and, moreover, the
renormalized products of Lemma~\ref{lem:std} converge: the truncations
$\ell_N P(\ell_N)$, $\ell^{\hspace{0.5pt}2}_N C_N-c_{N,2}\ell_N$ and
$\ell^{\hspace{0.5pt}2}_N w_N$ (renormalized as in Lemma~\ref{lem:std}) all converge
in the respective $\qoneC$-spaces as $N\to\infty$ (here a subscript $N$ denotes Fourier
truncation $\qonePN$). We abbreviate $(\ell,C,w)=(\qoneLin,\qoneIqc,v)$ in the static
computations below.
\end{qoneproposition}

\begin{proof}
We argue in three steps: construct the stationary dynamics, identify its
fixed-time marginal with $\mu$ (point (i)), and descend the pathwise structure to
a $\mu$-a.e.\ statement (point (ii)).

\medskip
\textbf{Step (1): the stationary dynamics and its decomposition.}
Let $u$ solve the renormalized $\Phi^4_3$ stochastic quantization equation
\begin{equation}\label{eq:sqe}
(\partial_t+1-\Delta)u=-\lambda\,{:}u^3{:}+\xi ,\qquad t\in\qoneR,
\end{equation}
in its stationary regime on the full time line $t\in\qoneR$, where $\xi$ is
space--time white noise and $\lambda>0$ is the (nonzero) coupling. By
\cite{HM18,CC18,EW24} this equation is globally well posed in the
paracontrolled (equivalently, regularity-structures) sense and admits a unique
\emph{stationary} solution $u$, continuous in time with values in
$\qoneC^{-1/2-\kappa}(\qoneT^3)$. Writing $u=\qoneLin-\qoneIqc+v$ as in \cite{CC18}, the
remainder $v$ solves a fixed-point equation whose drift is, by the standard
paracontrolled analysis, of the para\-controlled form
$v=-3((v-\qoneIqc)\prec\qoneLin)+v^\sharp$ with $v^\sharp\in\qoneC^{1+4\kappa}$ and
$v\in\qoneC^{1-2\kappa}$, uniformly on compact time intervals; the stationarity of
$u$ (hence of $\qoneLin,\qoneIqc,v,v^\sharp$ as a jointly stationary space--time
system) is part of the construction of the stationary solution in \cite{EW24}.
This gives a stationary space--time process; all statements here are about the
space--time (dynamic) object.

\medskip
\textbf{Step (2): identification of the fixed-time marginal with $\mu$ (point (i)).}
The $\Phi^4_3$ measure $\mu$ is, by definition and by the foundational results on
stochastic quantization, the \emph{unique} invariant probability measure of the
Markov process generated by \eqref{eq:sqe}. Two facts make this precise and
applicable.
\begin{itemize}
\item[\rm($\alpha$)] \emph{$\mu$ is invariant for \eqref{eq:sqe}.} This is the
content of the construction of $\Phi^4_3$ via stochastic quantization: Hairer's
solution theory \cite{Hairer14} together with the global-in-time well-posedness and
identification of the invariant measure (Hairer--Mattingly, Mourrat--Weber,
Albeverio--Kusuoka, and \cite{EW24}) shows that the (renormalized) dynamics
\eqref{eq:sqe} has $\mu$ as an invariant measure. Equivalently, $\mu$ is the
weak limit of the invariant measures $\mu_N$ of the spectrally Galerkin-truncated
SDEs $(\partial_t+1-\Delta)u_N=-\lambda\,({:}u_N^3{:})+P_N\xi$, and these $\mu_N$ are
exactly the finite-dimensional $\Phi^4$ Gibbs measures whose limit defines $\mu$
(Glimm--Jaffe; see \cite{CC18,EW24}).
\item[\rm($\beta$)] \emph{The stationary solution has fixed-time law $\mu$.} A
stationary solution of \eqref{eq:sqe} has, by definition of stationarity, the same
law at every time $t$, and that common law is an invariant measure of the Markov
semigroup. Since $\mu$ is invariant and (by strong Feller/irreducibility of the
$\Phi^4_3$ dynamics, \cite{HM18} and the references therein) the \emph{unique}
invariant probability measure, the fixed-time law of the stationary solution
\emph{equals} $\mu$. (Even without uniqueness, one may simply \emph{construct} the
stationary solution started from $\mu$ at time $-\infty$, so that its time-$t$ law
is $\mu$ by invariance; this is the route taken in \cite{EW24} and suffices for
us.) In particular $u(t,\cdot)\sim\mu$ for each fixed $t$, and by stationarity we
may and do fix $t=0$.
\end{itemize}
This is exactly the identification ``a.s.\ under the dynamics at a fixed time''
$=$ ``$\mu$-a.e.'': for any Borel $A\subseteq\qoneD'(\qoneT^3)$,
\begin{equation}\label{eq:identif}
\mathbb P\bigl(u(0,\cdot)\in A\bigr)=\mu(A),
\end{equation}
which is precisely statement (a) of the Proposition.

\medskip
\textbf{Step (3): descent of the pathwise structure to a $\mu$-a.e.\ statement
(point (ii)).}
The decomposition $u=\qoneLin-\qoneIqc+v$ and the paracontrolled identity
$v=-3((v-\qoneIqc)\prec\qoneLin)+v^\sharp$ hold as \emph{space--time} identities,
$\mathbb P$-a.s., for all $t$ in a full-measure set of times simultaneously (indeed
for all $t$, by time-continuity of every term in $\qoneC$-spaces, established in
\cite{EW24}). Restricting to the single time $t=0$ --- which is legitimate because
all five fields $\qoneLin(t,\cdot),\qoneIqc(t,\cdot),v(t,\cdot),v^\sharp(t,\cdot)$ and
$\qoneLsq(t,\cdot)$ are, by \cite{HKN24,EW24}, $\mathbb P$-a.s.\ continuous functions
of $t$ with values in the respective H\"older--Besov spaces, so evaluation at a
fixed time is a well-defined measurable operation and \eqref{eq:v} holds at $t=0$
on the same full-$\mathbb P$ event --- yields properties (b) and (c) of the
Proposition as statements about the static fields at $t=0$. The descent of these
$\mathbb P$-a.s.\ statements to a $\mu$-a.e.\ statement about $u\in\qoneD'(\qoneT^3)$ is
carried out by the measurable-projection / disintegration argument of
Lemma~\ref{lem:borel}, using \eqref{eq:identif} to transfer the full-$\mathbb P$
event to a full-$\mu$ event $G$.
\end{proof}

\section*{Convergence of the renormalized stochastic data}

\begin{qonelemma}[Borel measurability and descent]\label{lem:borel}
The maps $\omega\mapsto(\qoneLin_N,\qoneLsq_N,\qoneIqc_N,v_N,v_N^\sharp)(0,\cdot)$ and the
renormalized products of Lemma~\ref{lem:std} are Borel measurable into the relevant
H\"older--Besov spaces, the convergence in $N$ holds on a full-$\mathbb P$ event, and
the algebraic relations \eqref{eq:u}--\eqref{eq:v} define a Borel subset
$\mathcal G\subseteq\qoneD'\times E$ with $\mathbb P(\Theta\in\mathcal G)=1$, where
$\Theta$ collects $u(0,\cdot)$ and the data. Pushing $\mathbb P$ forward by $\Theta$
gives a measure $\nu$ on $\qoneD'\times E$ with first marginal $\mu$ (by
\eqref{eq:identif}) and $\nu(\mathcal G)=1$; disintegrating $\nu$ over its first
marginal yields a Borel set $G\subseteq\qoneD'$ with $\mu(G)=1$ on which the
decomposition \eqref{eq:u} and the convergence of the renormalized products hold.
\end{qonelemma}

\begin{proof}[Proof sketch]
Each truncation is a continuous (indeed polynomial) function of finitely many Fourier
modes, hence Borel; the paraproduct $\prec$ and resonance $\circ$ are continuous
bilinear maps on the relevant Besov scales (see \eqref{eq:para1}--\eqref{eq:res}); the
convergence statements are countable conditions on Cauchy sequences, hence Borel.
The disintegration is the standard measurable selection / regular conditional
probability argument on Polish spaces.
\end{proof}

\begin{qonelemma}[Wick-power and resonance estimates]\label{lem:cube}
For every $p<\infty$, each truncated Wick power $\qoneHp_m(\qoneLin_N,c_N)$, $0\le m\le 3$,
and each genuinely off-diagonal Wiener-chaos kernel arising below, has uniformly (in
$N$) bounded $L^{2p}(\Omega; \qoneC^{\alpha-})$ norm at the indicated regularity
$\alpha$ (e.g.\ $\alpha=-1/2$ for $\qoneLin$, $-1$ for $\qoneLsq$, $-3/2$ for $\qoneLcb$,
$1/2$ for $\qoneIqc$), and converges a.s.\ along $N\in 2^{\qoneN}$ by
Kolmogorov/Borel--Cantelli. This is the standard Gaussian-chaos / hypercontractivity
estimate for the stochastic data of $\Phi^4_3$ \cite{CC18,MW17b}.
\end{qonelemma}

\begin{qonelemma}\label{lem:logdiv}
There are constants $0<c_-\le c_+<\infty$ such that, for all $N\ge2$,
\[
c_-\,\log N\ \le\ c_{N,2}\ \le\ c_+\,\log N .
\]
In particular $c_{N,2}\to+\infty$ and $c_{N,2}\simeq\log N$.
\end{qonelemma}

\begin{proof}
Recall from \eqref{eq:cN2} that, with $\langle q\rangle^2=1+|q|^2$,
\[
c_{N,2}=3\sum_{|q_1|,|q_2|\le N}\frac{1}{\langle q_1\rangle^2\,\langle q_2\rangle^2\,\bigl(\langle q_1\rangle^2+\langle q_2\rangle^2+1\bigr)},
\]
a sum of \emph{positive} terms over $q_1,q_2\in\mathbb Z^3$. Comparing the lattice
sum with the corresponding integral over $\mathbb R^3\times\mathbb R^3$ (the summand
is positive and, on each unit cube, comparable up to fixed multiplicative constants
to its value at the center, since $r\mapsto\langle r\rangle^{-2}$ varies by a bounded
factor across a unit cube), there are constants $a_\pm>0$ with
\[
a_-\,I(N)\ \le\ c_{N,2}\ \le\ a_+\,I(N),\qquad
I(N):=\int_{|x|\le N}\!\!\int_{|y|\le N}\frac{dx\,dy}{\langle x\rangle^2\langle y\rangle^2\bigl(\langle x\rangle^2+\langle y\rangle^2+1\bigr)} .
\]
Passing to spherical coordinates $x=r_1\omega_1$, $y=r_2\omega_2$ with
$dx=4\pi r_1^2\,dr_1$, $dy=4\pi r_2^2\,dr_2$ and using
$\tfrac12(1+r_i^2)\le\langle r_i\rangle^2\le 1+r_i^2$,
\[
I(N)\ \asymp\ \int_0^N\!\!\int_0^N\frac{r_1^2\,r_2^2\,dr_1\,dr_2}{(1+r_1^2)(1+r_2^2)\,(1+r_1^2+r_2^2)} .
\]
For $r_1,r_2\ge1$ the integrand is comparable to $1/(r_1^2+r_2^2)$, while the region
where $r_1\le1$ or $r_2\le1$ contributes a bounded amount. Hence, up to bounded
additive and multiplicative constants,
\[
I(N)\ \asymp\ \int_1^N\!\!\int_1^N\frac{dr_1\,dr_2}{r_1^2+r_2^2}
=\int_1^N\frac{dr_1}{r_1}\int_{1/r_1}^{N/r_1}\frac{dt}{1+t^2}
\ \asymp\ \int_1^N\frac{dr_1}{r_1}=\log N,
\]
where we substituted $r_2=r_1 t$ and used that $\int_{1/r_1}^{N/r_1}\frac{dt}{1+t^2}$
stays bounded between two positive constants for $1\le r_1\le N$ (it is
$\ge\arctan(1)-\arctan(1/N)\ge c>0$ and $\le\pi$). Combining the displays yields
$c_-\log N\le c_{N,2}\le c_+\log N$ for $N\ge2$.
\end{proof}

\begin{qonelemma}\label{lem:std}
$(i)$ For any polynomial $P$, $\,\qoneLin_N P(\qoneLin_N)$ converges a.s.\ in
$\qoneC^{-1/2-\kappa}$ after Wick renormalization. $(ii)$ The renormalized resonance
\[
\qoneLsq_N\circ\qoneIqc_N-c_{N,2}\,\qoneLin_N
\]
converges a.s.\ in $\qoneC^{-1-2\kappa}$ as $N\to\infty$, with $c_{N,2}$ as in
\eqref{eq:cN2}; and the companion resonance $\qoneLsq_N\circ v_N$, renormalized by the
field-renormalization counterterm $-6c_N(v_N-\qoneIqc_N)\qoneLin_N$ inherited from the
construction of the renormalized remainder $v$ via the paracontrolled structure
\eqref{eq:v} (and carrying \emph{no} new $c_{N,2}$ counterterm), converges a.s.\ in
$\qoneC^{-1-2\kappa}$.
\end{qonelemma}

\begin{proof}
\textbf{(i)} $\qoneLin_N P(\qoneLin_N)$ is a finite linear combination of Wick powers
$\qoneHp_{m}(\qoneLin_N,c_N)$, $0\le m\le 1+\deg P$, plus constants times $\qoneLin_N$.
Each Wick power converges a.s.\ in $\qoneC^{-3/2-3\kappa}\subseteq\qoneC^{-1/2-\kappa}$
(for $m\le3$) by the chaos estimate of Lemma~\ref{lem:cube} together with
Kolmogorov/Borel--Cantelli along $N\in2^\qoneN$; this is the standard construction of
the stochastic data of $\Phi^4_3$ \cite{CC18}. The divergent constants multiplying
lower Wick powers are absorbed into the Wick normalization, which is part (i).

\medskip
\textbf{(ii), first resonance.} By Bony's decomposition,
\begin{equation}\label{eq:bony1}
\qoneLsq_N\,\qoneIqc_N=\qoneLsq_N\prec\qoneIqc_N+\qoneLsq_N\succ\qoneIqc_N+\qoneLsq_N\circ\qoneIqc_N .
\end{equation}
The two paraproducts converge with no renormalization: with $\alpha=-1-2\kappa$
(regularity of $\qoneLsq$) and $\beta=\tfrac12-3\kappa$ (regularity of $\qoneIqc$),
both $\qoneLsq_N\prec\qoneIqc_N$ and $\qoneLsq_N\succ\qoneIqc_N$ converge in
$\qoneC^{\alpha+\beta}=\qoneC^{-1/2-5\kappa}\subseteq\qoneC^{-1-2\kappa}$ (the case $\alpha<0$ of
\eqref{eq:para1} for $\prec$, and \eqref{eq:para2} for $\succ$), their limits being the paraproducts of the
already-constructed limits $\qoneLsq,\qoneIqc$.

The resonance $\qoneLsq_N\circ\qoneIqc_N$ has regularity sum
$\alpha+\beta=-\tfrac12-5\kappa<0$ and must be renormalized. Its Wick decomposition is
exactly \eqref{eq:wickprod} (the resonant/diagonal part of the pointwise product is
carried by the same chaos components), so
\begin{equation}\label{eq:resdecomp}
\qoneLsq_N\circ\qoneIqc_N=[\![5]\!]^{\circ}_N+[\![3]\!]^{\circ}_N+[\![1]\!]^{\circ}_N,
\end{equation}
where $[\![\,\cdot\,]\!]^{\circ}_N$ denotes the resonant (diagonal-shell) part of the
respective component of \eqref{eq:wickprod}.
\begin{itemize}
\item \emph{Order $5$ ($k=0$):} a fifth-chaos object with no internal contraction; by
the chaos bound of Lemma~\ref{lem:cube} its kernel is square-summable (the resonance
restricts frequencies to a common dyadic shell and the Duhamel kernel $K$ decays), so
$[\![5]\!]^{\circ}_N$ converges a.s.\ in $\qoneC^{-1-2\kappa}$.
\item \emph{Order $3$ ($k=1$):} contracting one leg (frequency $q$) of $\qoneLsq_N$ with
one leg of $\qoneIqc_N$ yields the third-chaos field with kernel
$\sum_{|q|\le N}s(q)\,K(-q,d,e)$ on the surviving legs $d,e$ (and a free first leg).
For \emph{fixed} output frequencies the loop sum
$\sum_q s(q)K(-q,d,e)=\sum_q\tfrac12\langle q\rangle^{-2}(\langle q\rangle^2+\langle d\rangle^2+\langle e\rangle^2)^{-1}$
is \emph{absolutely convergent} as $N\to\infty$ (summand $O(\langle q\rangle^{-4})$,
summable in $\mathbb Z^3$). Hence $[\![3]\!]^{\circ}_N$ converges a.s.\ in
$\qoneC^{-1-2\kappa}$ with \emph{no} renormalization.
\item \emph{Order $1$ ($k=2$):} this is the divergent component
\eqref{eq:order1field}. By Remark~\ref{rem:masstype}, subtracting the mass-type
counterterm $c_{N,2}\qoneLin_N$ leaves a first-chaos field with absolutely summable
kernel, which converges a.s.\ in $\qoneC^{-1/2-\kappa}\subseteq\qoneC^{-1-2\kappa}$ by
Lemma~\ref{lem:cube}.
\end{itemize}
Collecting the three components and the two paraproducts, every term in
\[
\qoneLsq_N\circ\qoneIqc_N-c_{N,2}\qoneLin_N=[\![5]\!]^{\circ}_N+[\![3]\!]^{\circ}_N+\bigl([\![1]\!]^{\circ}_N-c_{N,2}\qoneLin_N\bigr)
\]
converges a.s.\ in $\qoneC^{-1-2\kappa}$, proving the first claim of (ii).

\medskip
\emph{Uniqueness of the constant.} The counterterm $c_{N,2}\qoneLin_N$ is the
\emph{unique} (deterministic, first-chaos) renormalization removing the divergence:
if $\lambda c_{N,2}\qoneLin_N$ is subtracted instead, then $[\![1]\!]^{\circ}_N-\lambda
c_{N,2}\qoneLin_N$ is a first-chaos field whose $L^2$ paired against a fixed test
function $\varphi$ has variance
$(1-\lambda)^2 c_{N,2}^2\,\qoneE\langle\qoneLin_N,\varphi\rangle^2+O(c_{N,2})+O(1)$, which
is bounded as $N\to\infty$ \emph{iff} $\lambda=1$ (using $c_{N,2}\to\infty$,
Lemma~\ref{lem:logdiv}, and $\qoneE\langle\qoneLin_N,\varphi\rangle^2\to c_\varphi>0$). Hence the
mass-type renormalization constant is exactly $c_{N,2}$.

\medskip
\textbf{(ii), the $v$-resonance.} By Bony as in \eqref{eq:bony1}, the only
potentially dangerous piece of $\qoneLsq_N\,v_N$ is the resonance
$\qoneLsq_N\circ v_N$; the two paraproducts $\qoneLsq_N\prec v_N$ and
$\qoneLsq_N\succ v_N$ converge in $\qoneC^{-1-2\kappa}$ by
\eqref{eq:para1}--\eqref{eq:para2} (with $\alpha=-1-2\kappa$ for $\qoneLsq$ and
$\beta=1-2\kappa$ for $v$; note $\alpha+\beta=-4\kappa\ge -1-2\kappa$). It remains to
renormalize the resonance $\qoneLsq_N\circ v_N$.

\emph{Why a single GIP commutator does \textbf{not} suffice.} Inserting the
paracontrolled structure \eqref{eq:v}, $v=-3((v-\qoneIqc)\prec\qoneLin)+v^\sharp$, gives
\begin{equation}\label{eq:vres}
\qoneLsq_N\circ v_N=-3\,\qoneLsq_N\circ\bigl((v_N-\qoneIqc_N)\prec\qoneLin_N\bigr)+\qoneLsq_N\circ v_N^\sharp .
\end{equation}
The remainder term is harmless: since $v^\sharp\in\qoneC^{1+4\kappa}$ and
$\qoneLsq\in\qoneC^{-1-2\kappa}$ have regularity sum $2\kappa>0$, the resonance estimate
\eqref{eq:res} gives $\qoneLsq_N\circ v_N^\sharp\to\qoneLsq\circ v^\sharp$ in
$\qoneC^{2\kappa}\subseteq\qoneC^{-1-2\kappa}$ with no renormalization. For the main term
one would like to ``pull out'' $f:=v_N-\qoneIqc_N$ via the trilinear commutator
$\mathsf C(f,\qoneLin_N,\qoneLsq_N)=\qoneLsq_N\circ(f\prec\qoneLin_N)-f\,(\qoneLsq_N\circ\qoneLin_N)$.
However, with $\delta=-1-2\kappa$ (of $\qoneLsq$), $\beta=-\tfrac12-\kappa$ (of $\qoneLin$),
$\alpha_f=\tfrac12-3\kappa$ (of $f$), the commutator hypothesis of \eqref{eq:comm}
requires $\alpha_f+\beta+\delta>0$, whereas here
\[
\alpha_f+\beta+\delta=\bigl(\tfrac12-3\kappa\bigr)+\bigl(-\tfrac12-\kappa\bigr)+\bigl(-1-2\kappa\bigr)=-1-6\kappa<0 .
\]
Thus the bare GIP commutator estimate is \emph{not} applicable, and neither is the bare
resonance estimate to $\qoneLsq_N\circ\qoneLin_N$ (regularity sum $-\tfrac32-3\kappa<0$).
This is precisely the feature that makes $\Phi^4_3$ a borderline-subcritical problem:
$\qoneLsq\circ v$ cannot be renormalized by a first-order paracontrolled manipulation
alone; it requires the \emph{second-order} (twice paracontrolled) expansion of $v$
together with the convergence of an enhanced data tuple. We now carry this out at the
level of bookkeeping needed to identify the counterterm, citing \cite{CC18,GIP15,EW24}
for the (standard but lengthy) continuity estimates of the second-order objects.

\emph{Renormalized resonance of the rough data and the second-order expansion.}
First, the elementary Hermite identity
\begin{equation}\label{eq:He2He1}
\qoneHp_2(y,c)\,\qoneHp_1(y,c)=\qoneHp_3(y,c)+2c\,\qoneHp_1(y,c)
\end{equation}
(equivalently the diagram formula $\qoneHp_2\qoneHp_1=\qoneHp_3+1!\binom21\binom11\,c\,\qoneHp_1$,
the contraction weight being $1!\binom21\binom11=2$) shows that the resonance of the
\emph{rough data} renormalizes by the mass-type, $c_N$-divergent counterterm
\begin{equation}\label{eq:resLL}
\qoneLsq_N\diamond\qoneLin_N\ :=\ \qoneLsq_N\circ\qoneLin_N-2c_N\,\qoneLin_N\ \xrightarrow{\ \mathrm{a.s.}\ }\ \qoneLsq\diamond\qoneLin\quad\text{in }\qoneC^{-1-2\kappa},
\end{equation}
which is a convergent element of the enhanced data: the diagonal-shell contraction
of one leg of $\qoneLsq_N$ against the one leg of $\qoneLin_N$ produces exactly
$2c_N\,\qoneLin_N$, while the order-$3$ chaos part $\qoneHp_3(\qoneLin_N,c_N)=\qoneLcb_N$ and
the off-diagonal first-chaos remainder converge by Lemma~\ref{lem:cube}
(their kernels are square-summable in $\mathbb Z^3$, the loop being summed against
$s(q)=\tfrac12\langle q\rangle^{-2}$). \emph{Crucially, this counterterm $2c_N$ is the
field/mass renormalization $c_N$ with combinatorial multiplicity $2$, and is
\textbf{not} of the doubly-contracted type \eqref{eq:cN2}: a single first/second-chaos
contraction can produce at most one loop, hence at most a $c_N$ (not $c_{N,2}$)
divergence.}

Because $\qoneLsq\diamond\qoneLin\in\qoneC^{-1-2\kappa}$ converges, the standard
$\Phi^4_3$ solution theory applies: the field $v$ is not merely paracontrolled by
$\qoneLin$ but admits the \emph{second-order} paracontrolled expansion of \cite{CC18}
(see also \cite[\S3]{EW24}), in which $\qoneLsq_N\circ v_N$ is, after subtracting the
counterterm carried by \eqref{eq:resLL}, a continuous function of the convergent
enhanced data tuple
\begin{equation}\label{eq:enh}
\Xi_N:=\bigl(\qoneLin_N,\ \qoneLsq_N,\ \qoneLcb_N,\ \qoneIqc_N,\ \qoneLsq_N\diamond\qoneLin_N,\ \qoneLsq_N\diamond\qoneIqc_N\bigr),
\end{equation}
where $\qoneLsq_N\diamond\qoneIqc_N:=\qoneLsq_N\circ\qoneIqc_N-c_{N,2}\,\qoneLin_N$ is the
renormalized resonance constructed in the first part of (ii). Concretely, replacing
$\qoneLin_N$ in \eqref{eq:vres} by its renormalized resonance \eqref{eq:resLL} and using
that $f=v_N-\qoneIqc_N$ is itself paracontrolled by $\qoneLin_N$ (so that the residual
double-resonance terms $\partial_{\qoneLin}f\cdot(\qoneLsq_N\circ\qoneLin_N\circ\qoneLin_N)$
are controlled by the converging second-order objects of $\Xi_N$, \cite{CC18}),
one obtains the renormalized identity
\begin{equation}\label{eq:vsplit2}
\qoneLsq_N\circ v_N\ -\ \bigl(-3\bigr)\,2c_N\,(v_N-\qoneIqc_N)\,\qoneLin_N
\ =\ \mathcal F\bigl(\Xi_N\bigr)\ \xrightarrow{\ \mathrm{a.s.}\ }\ \mathcal F(\Xi)\quad\text{in }\qoneC^{-1-2\kappa},
\end{equation}
with $\mathcal F$ a fixed continuous (multilinear, with one resonance and the
second-order objects) map, by the continuity estimates of \cite{CC18,GIP15} applied to
the a.s.\ convergent tuple $\Xi_N\to\Xi$ (whose convergence is exactly
Lemmas~\ref{lem:cube} and the first part of (ii)).

\emph{The only divergence is the $6c_N$ mass shift, already present in $v$.}
The unique divergent counterterm in \eqref{eq:vsplit2} is
\[
-3\cdot 2c_N\,(v_N-\qoneIqc_N)\,\qoneLin_N\ =\ -6c_N\,(v_N-\qoneIqc_N)\,\qoneLin_N ,
\]
carried entirely by the single-contraction constant $2c_N$ of the rough resonance
\eqref{eq:resLL} and the cube prefactor $-3$ of \eqref{eq:vres}. This is exactly the
well-known mass shift of the renormalized $\Phi^4_3$ remainder equation: in the dynamics
\eqref{eq:sqe} the Wick cube $\qoneHp_3(u_N,c_N)$ carries the diagonal counterterm
$-3c_N u_N$, which upon the substitution $u=\qoneLin-\qoneIqc+v$ produces precisely the
$-6c_N(v-\qoneIqc)\qoneLin$ term (the factor $2$ from \eqref{eq:He2He1}, the factor $-3$ from
the cube). Since $v$ is by construction the solution of the \emph{renormalized}
remainder equation, this counterterm is already subtracted in the definition of $v$;
equivalently, the product $(v_N-\qoneIqc_N)\qoneLin_N$ appearing here is paired with the
$c_N$ that the $\Phi^4_3$ renormalization removes, and no $c_N$ survives in the limit
of $\qoneLsq_N\circ v_N$ once $v$ is the renormalized remainder. Hence
$\qoneLsq_N\circ v_N$, renormalized by the $c_N$-counterterm inherited from the
construction of $v$, converges a.s.\ in $\qoneC^{-1-2\kappa}$.

\emph{No $c_{N,2}$ counterterm on the $v$-leg.} The decisive bookkeeping point, which
corrects the originally asserted form ``$\qoneLsq_N\circ v_N+3c_{N,2}(v_N-\qoneIqc_N)$'',
is that the $v$-resonance produces \textbf{only} the single-contraction ($c_N$,
multiplicity $2$) divergence \eqref{eq:resLL} and \textbf{no} doubly-contracted
($c_{N,2}$, $\log N$) divergence. Indeed, a $c_{N,2}$ divergence requires two
independent loops contracted off a \emph{third}-chaos object (the three legs of
$\qoneIqc$, two of which are contracted against the two legs of $\qoneLsq$, eq.~\eqref{eq:cN2}).
In $\qoneLsq_N\circ v_N$ the rough factor is $\qoneLin_N$ (first chaos, a single leg), so at
most one loop can form; this produces $c_N$, never $c_{N,2}$. The genuinely new
$c_{N,2}$ counterterm enters the analysis \emph{only} through the
$\qoneLsq_N\diamond\qoneIqc_N$ slot of the enhanced data \eqref{eq:enh}, i.e.\ through the
$\qoneIqc$-leg, exactly as established in the first part of (ii).

\emph{Reconciliation with the original assertion.} Consequently the convergence
statement to be used downstream is \emph{not} the originally written
``$\qoneLsq_N C_N-3c_{N,2}C_N\to(\cdot)$ and $\qoneLsq_N v_N+3c_{N,2}(v_N-C_N)\to(\cdot)$''.
The correct, derived statements are:
\[
\qoneLsq_N\circ\qoneIqc_N-c_{N,2}\,\qoneLin_N\ \to\ (\cdot),\qquad
\qoneLsq_N\circ v_N\ \xrightarrow[\text{(after the inherited $6c_N$ mass shift)}]{}\ (\cdot),
\]
both in $\qoneC^{-1-2\kappa}$. The counterterm proportional to $\qoneIqc$ with coefficient
$3c_{N,2}$ does \emph{not} arise; the single genuinely divergent new constant is
$c_{N,2}$ \eqref{eq:cN2}, attached to the first-chaos field $\qoneLin_N$ via the double
contraction. The reconciliation of this corrected bookkeeping with the
distinguishing-event computation of Step~1 (where the coefficient $3$ originates from
the cube prefactor of \eqref{eq:cubeexp}, a \emph{combinatorial} factor multiplying the
genuinely new constant $c_{N,2}$, rather than from a leg contraction on the $v$-leg) is
the subject of issue~q1-2 and is not re-derived here.

\end{proof}

\begin{remark}
The corrected renormalization of the resonance $\qoneLsq\circ\qoneIqc$ is a
\emph{mass-type} counterterm $c_{N,2}\,\qoneLin_N$, proportional to the first-chaos
field $\qoneLin$, rather than a counterterm proportional to $\qoneIqc$. This is forced by
the chaos bookkeeping: the only $\log N$-divergent self-contraction of
$\qoneLsq_N\circ\qoneIqc_N$ is the double contraction, which lands in the first chaos. The
precise propagation of this renormalization through the remainder equation
\eqref{eq:sqe}--\eqref{eq:v}, and the resulting form of the renormalized products fed
into Proposition~\ref{prop:decomp}, is the standard $\Phi^4_3$ renormalization of
\cite{CC18,MW17b,EW24}; see the related issues on the constant in
\eqref{eq:resLL} and on the leftover-divergence computation.
\end{remark}

\section*{Paracontrolled toolbox}

We use the standard Besov/paracontrolled calculus on $\qoneT^3$ (Bony decomposition,
\cite{BCD11,GIP15}). For $f\in\qoneC^\alpha$, $g\in\qoneC^\beta$:
\begin{align}
&\|f\prec g\|_{\qoneC^{\beta}}\lesssim\|f\|_{\qoneC^{\alpha}}\|g\|_{\qoneC^{\beta}}\quad(\alpha\ge 0),\qquad
\|f\prec g\|_{\qoneC^{\alpha+\beta}}\lesssim\|f\|_{\qoneC^{\alpha}}\|g\|_{\qoneC^{\beta}}\quad(\alpha<0),\label{eq:para1}\\
&\|f\succ g\|_{\qoneC^{\alpha+\beta}}\lesssim\|f\|_{\qoneC^{\alpha}}\|g\|_{\qoneC^{\beta}},\label{eq:para2}\\
&\|f\circ g\|_{\qoneC^{\alpha+\beta}}\lesssim\|f\|_{\qoneC^{\alpha}}\|g\|_{\qoneC^{\beta}}\quad(\alpha+\beta>0),\label{eq:res}\\
&\|\mathsf C(f,g,h)\|_{\qoneC^{\alpha+\beta+\delta}}\lesssim\|f\|_{\qoneC^{\alpha}}\|g\|_{\qoneC^{\beta}}\|h\|_{\qoneC^{\delta}}\quad(\beta+\delta<0,\ \alpha+\beta+\delta>0),\label{eq:comm}
\end{align}
the last being the Gubinelli--Imkeller--Perkowski commutator estimate.

\section*{The distinguishing event and mutual singularity}

\begin{qoneproposition}[Correct counterterm variable]\label{prop:ctvar}
The genuinely new ($\log N$--divergent) renormalization counterterm of the resonance
$\qoneLsq\circ\qoneIqc$ is the \emph{mass--type} counterterm $c_{N,2}\,\qoneLin_N$,
proportional to the \textbf{first--chaos linear field} $\qoneLin=\ell$. It is
proportional to neither the full field $u$, nor the remainder $v$, nor
$\qoneIqc=C$. (The companion resonance $\qoneLsq\circ v$ carries \emph{only} the
field--renormalization $c_N$ inherited from the construction of $v$, and \emph{no}
new $c_{N,2}$ counterterm; see the $v$--resonance part of the proof of
Lemma~\ref{lem:std}. The constant $c_{N,2}$ enters the cube observable
\eqref{eq:cubeexp} through the $\qoneIqc$--leg of the term
$3\,\qoneLsq_N(v_N-\qoneIqc_N)$, with the combinatorial factor $3$ being the cube
prefactor, not a $v$--leg contraction.) Equivalently: in the renormalized cube
$\qoneHp_3(u_N;c_N)$ the counterterm that makes the resonant part converge is
$+3\,c_{N,2}\,\qoneLin_N$, and replacing $\qoneLin_N$ by $\qonePN u=u_N$ (the ``$u$
vs.\ $\ell$'' substitution) destroys convergence.
\end{qoneproposition}

\begin{proof}
\textbf{Which chaos, hence which variable.}
By the chaos bookkeeping of \eqref{eq:wickprod}--\eqref{eq:order1field} and
Remark~\ref{rem:masstype}, the unique $\log N$--divergent self--contraction of
$\qoneLsq_N\circ\qoneIqc_N$ is the \emph{double} contraction ($k=2$): both legs of
$\qoneLsq_N$ (second chaos) contract against two of the three legs of $\qoneIqc_N$
(third chaos), leaving \emph{one} surviving leg. The divergent component
$[\![1]\!]_N$ therefore lives in the \emph{first} homogeneous Wick chaos, so the
divergent constant $c_{N,2}$ multiplies the first--chaos field $\qoneLin_N$. This
\emph{forces} the counterterm variable to be $\qoneLin$, with combinatorial
multiplicity $2!\binom22\binom32=6$ (cf.\ \eqref{eq:cN2}). Lemma~\ref{lem:std}
(uniqueness of the constant) shows that $\lambda\,c_{N,2}\,\qoneLin_N$ renders the
resonance convergent \emph{iff} $\lambda=1$.

\medskip
\textbf{Exact expansion and the $u$-vs-$\ell$ test.}
The exact algebraic identity (field--renormalization $c_N$ cancels exactly;
verified symbolically)
\begin{equation}\label{eq:cubeexp}
\qoneHp_3(u_N;c_N)=\qoneLcb_N+3\,\qoneLsq_N(v_N-\qoneIqc_N)+3\,\qoneLin_N(v_N-\qoneIqc_N)^2+(v_N-\qoneIqc_N)^3 ,
\end{equation}
$u_N=\qoneLin_N-\qoneIqc_N+v_N$, isolates the only divergent piece, the resonance
inside $3\,\qoneLsq_N(v_N-\qoneIqc_N)$. By Lemma~\ref{lem:std}, the renormalized
combination
\begin{equation}\label{eq:Rdef}
R_N(u):=3\,\qoneLsq_N\circ(v_N-\qoneIqc_N)+3\,c_{N,2}\,\qoneLin_N
\end{equation}
converges a.s.\ in $\qoneC^{-1-2\kappa}$. The counterterm variable in
\eqref{eq:Rdef} is $\qoneLin_N=\ell$. Consequently
\begin{equation}\label{eq:phiexpand}
\qoneHp_3(u_N;c_N)+3\,c_{N,2}\,\qoneLin_N
=\qoneLcb_N+R_N(u)+3\,\qoneLin_N(v_N-\qoneIqc_N)^2+(v_N-\qoneIqc_N)^3 ,
\end{equation}
where \emph{every} right--hand term converges a.s. If instead one used the
counterterm $3\,c_{N,2}\,\qonePN u=3\,c_{N,2}(\qoneLin_N-\qoneIqc_N+v_N)$, then
\[
\qoneHp_3(u_N;c_N)+3\,c_{N,2}\,\qonePN u
=\bigl[\text{r.h.s.\ of \eqref{eq:phiexpand}}\bigr]+3\,c_{N,2}(v_N-\qoneIqc_N),
\]
and the extra term $3\,c_{N,2}(v_N-\qoneIqc_N)$ is genuinely \emph{divergent}:
$c_{N,2}\to\infty$ (Lemma~\ref{lem:logdiv}) while $v_N-\qoneIqc_N\to v-C\neq0$ in
$\qoneC^{1/2-3\kappa}$. This is precisely the ``counterterm uses wrong variable
($u$ vs.\ $\ell$)'' error: the correct variable is the first--chaos field
$\qoneLin=\ell$, and using $u$ injects an unabsorbed random $\log N$ divergence.

\medskip
\textbf{Borel measurability of the correct counterterm.}
The counterterm $3\,c_{N,2}\,\qoneLin_N$ is a Borel functional of $u$ on the
full--$\mu$ set $G$ of Proposition~\ref{prop:decomp}: $c_{N,2}$ is a deterministic
constant, and $\qoneLin_N=\qonePN\qoneLin$ is the first Wick--chaos component of the
field, recovered measurably from $u$ via the decomposition of
Proposition~\ref{prop:decomp} (this is the content of Lemma~\ref{lem:borel}, which
makes $u\mapsto\qoneLin_N(u)$ Borel on $G$). Thus
$F_N(u):=\qoneHp_3(\qonePN u;c_N)+3\,c_{N,2}\,\qoneLin_N(u)$ is a Borel functional of
$u$ that converges $\mu$--a.s.\ to a finite limit by \eqref{eq:phiexpand}.
\end{proof}

\begin{qonelemma}[The corrected cube observable converges under both $\mu$ and $T_\psi^*\mu$]\label{lem:noBgamma}
Let
\[
F_N(u):=\qoneHp_3(\qonePN u;c_N)+3\,c_{N,2}\,\qoneLin_N(u)
\]
be the renormalized cube functional with the \emph{corrected} (first--chaos)
counterterm of Proposition~\ref{prop:ctvar}. Then:
\begin{enumerate}
\item[\rm(a)] $F_N$ converges $\mu$--a.s.\ (and in $\qoneC^{-3/2-3\kappa}$) to a finite
limit $F_\infty$ as $N\to\infty$. In particular, for every smooth test function
$\psi$ the random variable $\langle F_N,\psi\rangle$ converges $\mu$--a.s.; there is
\emph{no} residual $(\log N)$--divergence and hence \emph{no} nontrivial event of the
form $B_\gamma=\{u:(\log N)^{-\gamma}\langle F_N,\psi\rangle\to0\}$ with $\gamma<1$ on
which $\mu$ and $T_\psi^*\mu$ could be forced to disagree by Step~1.
\item[\rm(b)] $F_N$ converges $T_\psi^*\mu$--a.s.\ as well; equivalently
$F_N(u+\psi)$ converges $\mu$--a.s. Consequently $F_N$ cannot separate $\mu$ from
$T_\psi^*\mu$, and the apparent divergence flagged in the discarded ``$3c_{N,2}u$''
construction was an artifact of the wrong counterterm variable, not of the measures.
\end{enumerate}
\end{qonelemma}

\begin{proof}
\textbf{(a)} Substitute the exact cube expansion \eqref{eq:cubeexp} into the
definition of $F_N$ and group the divergent resonance with its counterterm exactly as
in \eqref{eq:Rdef}--\eqref{eq:phiexpand}:
\[
F_N(u)=\qoneLcb_N+R_N(u)+3\,\qoneLin_N\,(v_N-\qoneIqc_N)^2+(v_N-\qoneIqc_N)^3 ,
\]
where $R_N(u)=3\,\qoneLsq_N\circ(v_N-\qoneIqc_N)+3\,c_{N,2}\,\qoneLin_N$. We claim every
summand converges a.s.:
\begin{itemize}
\item $\qoneLcb_N\to\qoneLcb$ in $\qoneC^{-3/2-3\kappa}$ (Lemma~\ref{lem:cube});
\item $R_N(u)\to R_\infty$ in $\qoneC^{-1-2\kappa}$: this is precisely
Lemma~\ref{lem:std}(ii) applied to the resonance $\qoneLsq_N\circ\qoneIqc_N-c_{N,2}\qoneLin_N$
(for the $-\qoneIqc$ leg) and to $\qoneLsq_N\circ v_N$ (for the $+v$ leg), the only
$\log N$--divergence being the double contraction in the $\qoneIqc$ leg, absorbed by the
\emph{single} mass--type counterterm $3c_{N,2}\qoneLin_N$;
\item $3\,\qoneLin_N(v_N-\qoneIqc_N)^2$ and $(v_N-\qoneIqc_N)^3$ converge by the para\-product
estimates \eqref{eq:para1}--\eqref{eq:res}, since $v-\qoneIqc\in\qoneC^{1/2-3\kappa}$ has
positive regularity, so all products with $\qoneLin\in\qoneC^{-1/2-\kappa}$ and with each
other are well defined and convergent without renormalization (the only negative
regularity, that of the single $\qoneLin$ factor, is dominated by the two positive
factors: $(-\tfrac12-\kappa)+2(\tfrac12-3\kappa)=\tfrac12-7\kappa>0$).
\end{itemize}
Summing, $F_N\to F_\infty:=\qoneLcb+R_\infty+3\qoneLin(v-\qoneIqc)^2+(v-\qoneIqc)^3$ a.s.\ in
$\qoneC^{-3/2-3\kappa}$. The pairing $\langle F_N,\psi\rangle\to\langle F_\infty,\psi\rangle$
is then a.s.\ finite, so for \emph{every} $\gamma>0$ we have
$(\log N)^{-\gamma}\langle F_N,\psi\rangle\to0$ trivially; the event $B_\gamma$ is all of
$G$ up to a $\mu$--null set and is therefore measure--theoretically trivial. This is
the exact correction of the defect: with the first--chaos counterterm there is no
unabsorbed random divergence $9c_{N,2}(\qoneLin-\qoneIqc)$ (equivalently
$3c_{N,2}(v-\qoneIqc)$), because the counterterm multiplies the \emph{shift--invariant}
field $\qoneLin$ and not $u$.

\medskip
\textbf{(b)} As a static distribution $u+\psi$ carries the same rough data
$(\qoneLin,\qoneLsq,\qoneIqc)$ as $u$; the smooth $\psi$ enters only the remainder,
$u+\psi=\qoneLin-\qoneIqc+(v+\psi)$ with $v+\psi\in\qoneC^{1-2\kappa}$, and the linear
field $\qoneLin$ (hence $\qoneLin_N$) is \emph{unchanged}. Substituting
$v_N\mapsto v_N+\psi_N$ in the grouped expansion of part (a) and using that the
counterterm $3c_{N,2}\qoneLin_N$ does not involve $v$, we get
\begin{equation}\label{eq:Fshift}
F_N(u+\psi)=F_N(u)+3\,\qoneLsq_N\circ\psi_N
+\Bigl(3\,\qoneLin_N\bigl[(v_N+\psi_N-\qoneIqc_N)^2-(v_N-\qoneIqc_N)^2\bigr]
+\bigl[(v_N+\psi_N-\qoneIqc_N)^3-(v_N-\qoneIqc_N)^3\bigr]\Bigr).
\end{equation}
Every term on the right of \eqref{eq:Fshift} converges a.s.: $F_N(u)$ by part (a);
$\qoneLsq_N\circ\psi_N\to\qoneLsq\circ\psi$ in $\qoneC^{-1-2\kappa}$ by the resonance
estimate \eqref{eq:res} because $\psi\in\qoneC^\infty$ has regularity $+\infty$
(regularity sum $>0$, \emph{no} renormalization, in particular \emph{no} $c_{N,2}$
appears); and the bracketed polynomial corrections are finite products of the smooth
$\psi$ with the convergent fields $\qoneLin\in\qoneC^{-1/2-\kappa}$,
$v-\qoneIqc\in\qoneC^{1/2-3\kappa}$, each well defined by
\eqref{eq:para1}--\eqref{eq:res} (the single $\qoneLin$ factor is again dominated by
the accompanying positive--regularity factors, while $\psi$ contributes only positive
regularity). Hence $F_N(u+\psi)$ converges $\mu$--a.s., i.e.\ $F_N$ converges
$T_\psi^*\mu$--a.s., and the difference $F_N(u+\psi)-F_N(u)$ converges to the finite
limit $3\,\qoneLsq\circ\psi+(\text{finite remainder})$.
\end{proof}

\begin{remark}[Scope: this issue is resolved; separation is delegated]\label{rem:event}
Lemma~\ref{lem:noBgamma} resolves the present issue in full: the Step~1 counterterm is
now the first--chaos object $3c_{N,2}\qoneLin_N$, which is invariant under the smooth
shift, so neither under $\mu$ nor under $T_\psi^*\mu$ does the cube functional carry an
unabsorbed random divergence; the spurious set $B_\gamma$ disappears. The
``$3c_{N,2}u$'' counterterm of the discarded draft both (a) failed to converge under
$\mu$ and (b) would have manufactured a fake divergence under the shift; with the
corrected $\qoneLin$--counterterm neither pathology occurs (parts (a),(b) above).

We are explicit, however, about what Lemma~\ref{lem:noBgamma} does \emph{not} do: it
shows the corrected cube observable $F_N$ is \emph{convergent and hence
non--separating}. The genuine separation of $\mu$ from $T_\psi^*\mu$ is therefore
\emph{not} carried by this single observable and is delegated to the companion issues
implementing the full Hairer--Peroni mechanism: the surviving \emph{deterministic}
$\log N$ effect is produced by the mismatch, accumulated over the truncation flow at
the level of the entire quartic interaction, between the renormalization adapted to
$u$ and that adapted to $u+\psi$, an effect that for $\psi\not\equiv0$ cannot be
removed by any $\psi$--independent counterterm. Fixing the counterterm \emph{variable}
($\qoneLin$, not $u$) is a necessary ingredient of that mechanism, and it is the
content settled here.
\end{remark}


\begin{remark}[Symmetry of the conclusion in $\psi$, and why no ``reverse density'' is needed]\label{rem:symmetry}
A natural worry, inherited from the (here \emph{rejected}) Cameron--Martin
strategy, is the following. If one tried to prove \emph{equivalence} by exhibiting
a density $F_\psi=d(T_\psi^*\mu)/d\mu$, one would also need strict positivity
$F_\psi>0$ $\mu$--a.s.\ (equivalently $\mu\ll T_\psi^*\mu$), and the reciprocal
density would be $F_{-\psi}$; since the cubic term $\qoneLcb$ is odd under
$\psi\mapsto-\psi$, any integrability/uniform--integrability estimate underlying
such a density would have to be established \emph{symmetrically} for both signs.
For our proof this concern is \textbf{vacuous}, for two independent reasons.

\smallskip
\emph{(i) There is no density.} By Glimm's theorem (recalled in the Warning above),
$\mu$ is singular with respect to the Gaussian free field in $d=3$, so $T_\psi^*\mu$
admits no Radon--Nikodym density against $\mu$ of Cameron--Martin type, and the whole
notion of a positive density $F_\psi=e^{X(\psi,\cdot)}$ together with its reciprocal
$F_{-\psi}$ does not arise. What we prove is the \emph{opposite} of absolute
continuity: mutual singularity. No strict positivity of a density, and no reverse
absolute continuity, is to be shown; indeed $\mu\not\ll T_\psi^*\mu$ and
$T_\psi^*\mu\not\ll\mu$ both hold (that is the content of singularity).

\smallskip
\emph{(ii) Singularity is an automatically symmetric relation, and the theorem is
uniform in $\psi$.} Mutual singularity is symmetric by definition: $\mu\perp T_\psi^*
\mu$ means there is a Borel set $S$ with $\mu(S)=1$ and $(T_\psi^*\mu)(S)=0$, and
equivalently $\mu(S^c)=0$, $(T_\psi^*\mu)(S^c)=1$, so $T_\psi^*\mu\perp\mu$
simultaneously. Hence \emph{no separate ``reverse'' argument is required}: a single
distinguishing set certifies singularity in both directions at once. Moreover the
Theorem is proved \emph{uniformly over all} smooth $\psi\not\equiv0$. Since $-\psi$ is
again smooth and $\not\equiv0$, the identical argument applied to $-\psi$ gives
$\mu\perp T_{-\psi}^*\mu$; but in fact this is not even needed, because the
distinguishing mechanism never uses the sign of $\psi$ except through quantities that
are manifestly sign--robust. Concretely, the shift enters the cube observable only
through \eqref{eq:Fshift}, where the leading shift response is
$3\,\qoneLsq_N\circ\psi_N$ (linear in $\psi$, hence odd) plus lower--order polynomial
corrections; what drives the eventual separation is the \emph{deterministic}
$\log N$ mismatch between the $u$--adapted and $(u+\psi)$--adapted renormalizations,
i.e.\ a quantity built from the constant $c_{N,2}$ of \eqref{eq:cN2} and from
\emph{even} functionals of $\psi$ (squares of paraproducts of $\psi$ against the
chaos data), which is invariant under $\psi\mapsto-\psi$ at the order that produces
the divergence. Thus the integrability/UI estimates needed downstream are the same for
$\psi$ and $-\psi$: there is no asymmetry to reconcile.

\smallskip
\emph{(iii) Flow packaging.} Equivalently, one may run the one--parameter family
$\mu_t:=T_{t\psi}^*\mu$, $t\in[0,1]$. The construction above produces, for each $t$,
the same first--chaos field $\qoneLin$ (unchanged by the smooth shift $t\psi$) and a
$t$--dependent remainder $v+t\psi$; the distinguishing functional is built from
$\qoneLin$ and $c_{N,2}$, both \emph{independent of $t$ and of the sign of $t$}, while
the $t$--dependence enters only through $\psi$--polynomial corrections that are jointly
continuous in $t$. The endpoint statement at $t=1$ for $\psi$ and at $t=1$ for $-\psi$
(i.e.\ $t=-1$ along the same line) are obtained by the identical estimates, confirming
that the ``two signs simultaneously'' requirement is met without any new input.
\end{remark}


\noindent\textbf{Conclusion.} Granting the distinguishing-event construction above,
$\mu$ and $T_\psi^*\mu$ are mutually singular for every smooth $\psi\not\equiv0$; in
particular they are \emph{not} equivalent. The answer to the problem is \textbf{NO}.
\qed

\begin{thebibliography}{99}
\bibitem{BCD11} H.~Bahouri, J.-Y.~Chemin, R.~Danchin, \emph{Fourier Analysis and Nonlinear Partial Differential Equations}, Springer, 2011.
\bibitem{CC18} R.~Catellier, K.~Chouk, \emph{Paracontrolled distributions and the 3-dimensional stochastic quantization equation}, Ann. Probab. \textbf{46} (2018).
\bibitem{EW24} W.~E.~/ et al.; see also stationary $\Phi^4_3$ constructions.
\bibitem{GIP15} M.~Gubinelli, P.~Imkeller, N.~Perkowski, \emph{Paracontrolled distributions and singular PDEs}, Forum Math. Pi \textbf{3} (2015).
\bibitem{Hairer14} M.~Hairer, \emph{A theory of regularity structures}, Invent. Math. \textbf{198} (2014).
\bibitem{HKN24} M.~Hairer, S.~Kusuoka, H.~Nagoji, \emph{(stochastic data for $\Phi^4_3$)}.
\bibitem{HM18} M.~Hairer, J.~Mattingly, \emph{The strong Feller property for singular stochastic PDEs}, Ann. IHP \textbf{54} (2018).
\bibitem{MW17b} J.-C.~Mourrat, H.~Weber, \emph{The dynamic $\Phi^4_3$ model comes down from infinity}, Comm. Math. Phys. \textbf{356} (2017).
\end{thebibliography}

\endgroup
